\documentclass[11pt]{article}

\usepackage[margin=1in]{geometry}
\usepackage{amsmath, amssymb}
\usepackage[T1]{fontenc}
\usepackage[utf8]{inputenc}
\usepackage{lmodern}
\usepackage{enumitem}
\usepackage{hyperref}

\setlist[itemize]{topsep=4pt, itemsep=2pt, parsep=0pt}

\title{6.8610 Spring 2026 Lecture 6 Notes: Pretraining and Fine-Tuning}
\author{Junru Ren with help from Codex}
\date{February 24, 2026}

\begin{document}
\maketitle

\section{Pretraining and Fine-Tuning}

\subsection{Motivation}
The lecture studies the modern transfer-learning pipeline for NLP: instead of training a system from scratch on a small task-specific dataset, first train a large general-purpose model and then adapt it to downstream tasks with much less specialized data.

\begin{itemize}
    \item Training a customer-service or task-oriented chatbot from scratch would require a dataset that teaches the model \emph{everything}: tokenization, syntax, conversation structure, politeness, domain knowledge, and task behavior.
    \item Collecting that much task-specific data is expensive and slow, especially when it requires skilled annotators.
    \item Pretraining solves this by learning general language and world knowledge from a large corpus, while fine-tuning communicates the task-specific behavior we actually want.
    \item Once pretraining enters the picture, the \textbf{optimization path} matters more. We are no longer just asking for any minimum of the loss; we want a basin from which small downstream updates generalize instead of merely memorizing.
    \item In practice, useful systems often require multiple adaptation stages rather than a single train-then-deploy step.
\end{itemize}

\subsection{Key Concepts}
\textbf{Definition: Pretraining}

Pretraining is the process of optimizing a model on a large, general-purpose corpus so that it learns broadly useful representations of language and related structure before any specific downstream task is imposed.

\textbf{Definition: Fine-Tuning}

Fine-tuning is the process of continuing training from a pretrained initialization on a smaller task- or domain-specific dataset in order to specialize the model's behavior.

\begin{itemize}
    \item \textbf{Base model:} a model pretrained on broad text that can model language fluently but is not reliably instruction-following or question-answering useful.
    \item \textbf{Instruction tuning:} fine-tuning a base model so that it responds usefully to natural-language prompts, questions, or task instructions.
    \item \textbf{In-context learning:} the ability to infer a task from examples or input-output pairs placed directly in the prompt; this can emerge from pretraining but can also be explicitly strengthened.
    \item \textbf{Refusal / unlearning:} a major goal of post-pretraining adaptation is to disable undesirable outputs, though in practice this often behaves more like \emph{refusal} than genuine forgetting.
    \item \textbf{Continued pretraining:} keep training on in-domain unlabeled text before or alongside downstream supervised fine-tuning.
    \item \textbf{Synthetic fine-tuning data:} use a stronger model to generate QA pairs, instructions, or additional domain documents when human-labeled data is scarce.
    \item \textbf{Data mixing:} pretraining corpora are engineered mixtures, not just arbitrary web scrapes.
    \item \textbf{Quality filtering:} remove low-quality or harmful text and keep samples that resemble trusted sources.
    \item \textbf{Parameter-efficient fine-tuning (PEFT):} adapt only a small subset of model parameters, such as bias terms or layer-normalization parameters, while leaving the rest frozen.
    \item \textbf{Shallow fine-tuning:} many downstream updates appear to surface or reweight capabilities already learned in pretraining rather than build new capabilities from scratch.
\end{itemize}

\subsection{Core Models and Equations}
\textbf{Transfer-learning pipeline.} Let \(D_{\mathrm{PT}}\) be a pretraining corpus and \(D_{\mathrm{FT}}\) a downstream dataset. A standard abstraction is
\[
\theta_{\mathrm{PT}}
=
\arg\max_{\theta}\sum_{(x,y)\in D_{\mathrm{PT}}}\log p_{\theta}(y \mid x),
\]
followed by
\[
\theta_{\mathrm{FT}}
=
\arg\max_{\theta}\sum_{(x,y)\in D_{\mathrm{FT}}}\log p_{\theta}(y \mid x)
\quad\text{initialized at } \theta=\theta_{\mathrm{PT}}.
\]
The lecture explicitly emphasized that the initialization matters because the optimizer is now navigating a highly nonconvex landscape.

\textbf{Gradient-descent view.} The slides wrote this in update form:
\[
\theta' \leftarrow \theta + \nabla_{\theta}\log p(y_{\mathrm{PT}} \mid x_{\mathrm{PT}}),
\]
\[
\theta'' \leftarrow \theta' + \nabla_{\theta'}\log p(y_{\mathrm{FT}} \mid x_{\mathrm{FT}}).
\]
The point is not the exact optimizer, but that the downstream update starts from a non-random parameter setting.

\textbf{Common pretraining objectives.}
\begin{itemize}
    \item \textbf{Autoregressive next-token prediction:}
    \[
    p_{\theta}(x_{1:T}) = \prod_{t=1}^{T} p_{\theta}(x_t \mid x_{1:t-1}).
    \]
    \item \textbf{Masked language modeling:}
    \[
    \max_{\theta} \sum_{t \in \mathcal{M}} \log p_{\theta}(x_t \mid x_{\setminus \mathcal{M}}),
    \]
    where \(\mathcal{M}\) is the set of masked positions.
    \item \textbf{Infilling / fill-in-the-middle:}
    \[
    \max_{\theta}\log p_{\theta}(x_{\mathrm{missing}} \mid x_{\mathrm{prefix}}, x_{\mathrm{suffix}}).
    \]
    \item \textbf{Ranking / discrimination:}
    \[
    p_{\theta}(y=1 \mid s_1, s_2),
    \]
    where the model predicts whether a sentence or sentence pair is real, ordered correctly, or preferred over a corrupted alternative.
\end{itemize}

\textbf{Importance weighting for data mixing.} If \(\hat{p}_{\mathrm{PT}}\) is a language model fit to pretraining data and \(\hat{p}_{\mathrm{FT}}\) is fit to downstream data, one can weight each pretraining example \(x\) by
\[
w(x) = \frac{\hat{p}_{\mathrm{FT}}(x)}{\hat{p}_{\mathrm{PT}}(x)}.
\]
Then the effective resampled distribution is
\[
p_{\mathrm{PT}}(x)\, w(x)
=
p_{\mathrm{PT}}(x)\frac{\hat{p}_{\mathrm{FT}}(x)}{\hat{p}_{\mathrm{PT}}(x)}
\approx
p_{\mathrm{FT}}(x),
\]
which upweights pretraining samples that look more like the eventual target domain.

\textbf{Adversarial reweighting intuition.} In methods such as DoReMi, one first trains a small reference model to identify hard subpopulations and then trains a larger model on an adversarially reweighted mixture. A representative form is
\[
\min_{\theta}\max_{q \in \Delta^K}
\sum_{k=1}^{K} q_k \big(\mathcal{L}_k(\theta) - \mathcal{L}_k(\theta_{\mathrm{ref}})\big),
\]
where \(q\) concentrates weight on subsets that remain difficult.

\textbf{Scaling laws.} The lecture did not dwell on formulas, but the core idea is that empirical loss obeys power-law behavior as a function of model size and dataset size. A representative abstraction is
\[
\mathcal{L}(N, D) \approx aN^{-\alpha} + bD^{-\beta} + c,
\]
where \(N\) is model size and \(D\) is dataset size.

\subsection{Algorithms / Architectures}
\textbf{1. Practical product pipeline}
\begin{itemize}
    \item Pretrain on large-scale text data.
    \item Instruction-tune on many tasks so the model becomes useful as a general assistant.
    \item Continue pretraining on in-domain text if the downstream domain is specialized.
    \item Optionally generate synthetic in-domain data.
    \item Fine-tune on human-written downstream examples.
\end{itemize}

\textbf{2. Choosing a pretraining objective}
\begin{itemize}
    \item Use autoregressive modeling when generation is central.
    \item Use masked language modeling when bidirectional representations are valuable.
    \item Use infilling when the deployment task resembles editing or code completion.
    \item Use ranking or corruption-detection objectives when the goal is to learn transferable structure without full generation.
\end{itemize}

\textbf{3. Designing the pretraining corpus}
\begin{itemize}
    \item Start from sources that are legally usable and broad enough to teach general capabilities.
    \item Mix sources intentionally rather than opportunistically.
    \item Filter for quality using trusted high-quality corpora as positive examples.
    \item Include code, metadata around code, and multilingual text when those capabilities are important.
    \item Use scaling-law heuristics to balance model size and corpus size.
\end{itemize}

\textbf{4. Continued pretraining and downstream adaptation}
\begin{itemize}
    \item Continue language-model training on in-domain unlabeled text.
    \item Fine-tune only on the outputs the model must generate at deployment time.
    \item If labeled examples are scarce, synthesize instruction or QA pairs from domain documents.
    \item If domain coverage is too narrow, synthesize additional plausible domain documents from the real corpus.
\end{itemize}

\textbf{5. Fine-tuning data selection}
\begin{itemize}
    \item Use influence-style methods to identify examples that most affect validation performance.
    \item Alternatively, hand-curate a very small but high-quality set of examples.
    \item Prefer quality over sheer volume when the base model is already strong.
\end{itemize}

\textbf{6. What parameters to update}
\begin{itemize}
    \item Update only embeddings in early transfer-learning setups.
    \item Freeze the base model and train only conversion layers or a classifier head.
    \item In modern PEFT, update only globally influential parameters such as bias terms or layer-normalization parameters.
    \item In full fine-tuning, update all parameters, but the resulting behavioral changes can still remain shallow.
\end{itemize}

\subsection{Important Intuitions from the Professor}
\begin{itemize}
    \item \textbf{From-scratch training on task data is usually the wrong baseline.} The lecture repeatedly stressed that a downstream dataset should not be responsible for teaching both the task and the whole structure of language.
    \item \textbf{Pretraining is partly about optimization geometry.} It is not just ``more data first.'' The initializer places the optimizer in a region where downstream updates can generalize.
    \item \textbf{A base model is not a product.} Pretraining on internet text yields a model of language, but not necessarily a model that follows instructions or answers questions in useful ways.
    \item \textbf{Code teaches more than code.} Training on code can improve reasoning and state-tracking behavior even on non-code tasks.
    \item \textbf{Multilingual pretraining induces shared latent structure.} The lecture highlighted evidence that models learn language-independent abstractions such as noun/verb/adjective structure.
    \item \textbf{Many modern improvements come from data design, not just model architecture.} Corpora must be mixed, filtered, and licensed intentionally.
    \item \textbf{Fine-tuning often surfaces existing mechanisms rather than teaching new ones.} If a thousand examples or a tiny parameter subset can dramatically change behavior, that suggests the core machinery was already learned during pretraining.
    \item \textbf{Refusal is fragile and not the same as forgetting.} Safety behavior can often be undone by tiny amounts of new fine-tuning or even by benign-seeming distributional shifts.
    \item \textbf{Trying to force new knowledge into a model can backfire.} If the data is thin and inconsistent, the model may instead learn that guessing or regurgitating is acceptable, which can raise hallucination rates.
\end{itemize}

\subsection{Examples}
\begin{itemize}
    \item \textbf{Customer-service chatbot:} the motivating example throughout the lecture was a website assistant that needs politeness, style, task following, and company-specific knowledge.
    \item \textbf{Word2vec-era transfer:} early systems pretrained embeddings and then froze them as initialization for downstream models.
    \item \textbf{Infilling for code assistants:} delete a function or middle span from a codebase and train the model to reconstruct it from prefix and suffix context.
    \item \textbf{BooksCorpus controversy:} BERT-era data pipelines relied on scraped ebooks whose permissions later became contentious, making some influential results hard to reproduce exactly.
    \item \textbf{The Pile:} an open, permissively licensed, terabyte-scale corpus that contrasts with opaque corporate pretraining mixtures.
    \item \textbf{Multilingual latent alignment:} multilingual models exhibit clustering by linguistic category rather than by language identity alone.
    \item \textbf{Reversal Curse:} fine-tuning can teach a fact in one direction, such as ``Daphne Barrington directed \emph{A Journey Through Time},'' without teaching the reversed query reliably.
    \item \textbf{Unexpected cross-task transfer:} training on buggy code can induce bizarre behavioral changes, including shifts in the model's apparent personality.
    \item \textbf{Hallucination from overteaching:} pushing too hard on fine-tuning for facts the model did not deeply know can reduce held-out accuracy and encourage unsupported guesses.
\end{itemize}

\subsection{Common Pitfalls / Limitations}
\begin{itemize}
    \item \textbf{Data licensing and reproducibility problems:} if the training mixture cannot legally be released, the scientific result may become difficult or impossible to reproduce exactly.
    \item \textbf{Poor corpus design:} indiscriminately scraping text can waste compute on low-quality, stale, or harmful data.
    \item \textbf{Synthetic data limits:} synthetic pretraining may not teach world knowledge, social structure, or the nuanced pragmatics of real language.
    \item \textbf{Temporal staleness:} both natural-language facts and codebases change over time, so pretraining alone can leave a model outdated.
    \item \textbf{Fine-tuning brittleness:} safety and refusal behaviors can be reversed by small downstream changes.
    \item \textbf{Logical incompleteness:} learning one direction of a relation does not guarantee learning its reverse or broader implications.
    \item \textbf{Unexpected side effects:} fine-tuning on one domain can change unrelated behaviors.
    \item \textbf{Hallucination risk:} trying to teach new facts with insufficiently rich data can make the model more willing to invent answers.
\end{itemize}

\subsection{Connections to Other Lectures}
\begin{itemize}
    \item This lecture ties back to Lecture 2, where word2vec already introduced the idea of learning useful representations from one task and reusing them elsewhere.
    \item It builds directly on Lecture 3 and Lecture 4, because modern pretraining and fine-tuning are overwhelmingly implemented with sequence models and especially Transformers.
    \item The lecture revisits BERT from Lecture 4 as an example of a pretraining objective designed for transferable representations rather than direct generation.
    \item It also connects to Lecture 5 on evaluation: pretraining corpora, continued pretraining, instruction tuning, and safety fine-tuning all require careful validation of what behavior actually changed.
    \item Conceptually, this lecture marks the shift from \emph{how to parameterize a model} to \emph{how to use staged training and data design to make that model useful}.
\end{itemize}

\subsection{Exam-Relevant Summary}
\begin{itemize}
    \item Pretraining learns general-purpose language capabilities; fine-tuning adapts those capabilities to specific tasks or domains.
    \item The practical product pipeline is often:
    \[
    \text{pretrain} \rightarrow \text{instruction tune} \rightarrow \text{in-domain continued pretrain} \rightarrow \text{task-specific fine-tune}.
    \]
    \item Important pretraining objectives include next-token prediction, masked language modeling, infilling, and ranking/corruption detection.
    \item Importance weighting for domain adaptation uses
    \[
    w(x)=\frac{\hat{p}_{\mathrm{FT}}(x)}{\hat{p}_{\mathrm{PT}}(x)}
    \]
    to upweight pretraining samples that resemble the downstream domain.
    \item Data mixtures should be designed, filtered, and legally usable; they are not merely arbitrary internet dumps.
    \item Training on code can improve non-code reasoning; multilingual pretraining can induce shared cross-lingual representations.
    \item Continued pretraining on in-domain unlabeled text often gives consistent gains before final supervised fine-tuning.
    \item Very small, carefully curated instruction datasets can rival much larger noisy ones.
    \item Parameter-efficient fine-tuning can update only a tiny subset of parameters and still strongly change behavior.
    \item Fine-tuning is often shallow: it mostly surfaces or steers capabilities already learned in pretraining.
    \item Refusal and unlearning are fragile; they usually do not erase underlying knowledge.
    \item Overaggressive fact fine-tuning can increase hallucination by teaching the model that unsupported guessing is acceptable.
\end{itemize}

\end{document}
